\section{Disparadores del Subsistema de Publicidad (Ismael)}

Para garantizar la integridad de los datos y cumplir con los requisitos de auditoría, se han implementado tres disparadores (triggers) en el SGBD. Estos mecanismos permiten automatizar validaciones de negocio complejas y mantener un histórico de modificaciones críticas.

\subsection{Trigger: Validación de Coherencia Temporal}
\textbf{Motivación:} Este disparador implementa una regla de negocio fundamental (\textit{Business Rule}): un anuncio no puede tener una fecha de finalización anterior a la fecha actual. Esto evita errores humanos al programar campañas o cerrar anuncios erróneamente en el pasado.

\textbf{Descripción:} Se ejecuta antes de insertar o actualizar (\texttt{BEFORE INSERT OR UPDATE}) en la tabla \texttt{ANUNCIO}. Compara la nueva fecha de fin con la fecha del sistema truncada (sin hora).

\begin{lstlisting}[language=SQL, caption={Trigger TRG\_PUB\_VALIDAR\_FECHAS}]
CREATE OR REPLACE TRIGGER TRG_PUB_VALIDAR_FECHAS
BEFORE INSERT OR UPDATE ON ANUNCIO
FOR EACH ROW
BEGIN
    -- Validamos contra la fecha actual (SYSDATE).
    -- Regla: No se puede crear o modificar un anuncio para que termine en el pasado.
    IF :NEW.FECHAFIN < TRUNC(SYSDATE) THEN
        RAISE_APPLICATION_ERROR(-20001, 
            'Error de Negocio: La fecha de fin del anuncio no puede ser anterior a la fecha actual.');
    END IF;
END;
/
\end{lstlisting}

\subsection{Trigger: Auditoría de Cambios en Vigencia}
\textbf{Motivación:} La modificación de la fecha de finalización de un anuncio tiene implicaciones contractuales y económicas. Por ello, es obligatorio registrar cualquier cambio en este campo para futuras auditorías o reclamaciones.

\textbf{Descripción:} Este disparador salta únicamente cuando se actualiza la columna \texttt{FECHAFIN}. Guarda automáticamente el estado anterior y el nuevo en la tabla histórica \texttt{HISTORIAL\_CAMBIOS\_FECHA}.

\begin{lstlisting}[language=SQL, caption={Trigger TRG\_AUDITAR\_CAMBIO\_FECHA}]
CREATE OR REPLACE TRIGGER TRG_AUDITAR_CAMBIO_FECHA
AFTER UPDATE OF FECHAFIN ON ANUNCIO
FOR EACH ROW
BEGIN
    INSERT INTO HISTORIAL_CAMBIOS_FECHA (IDANUNCIO, FECHA_ANTIGUA, FECHA_NUEVA)
    VALUES (:OLD.IDANUNCIO, :OLD.FECHAFIN, :NEW.FECHAFIN);
END;
/
\end{lstlisting}

\subsection{Trigger: Validación Semántica de Características}
\textbf{Motivación:} La tabla \texttt{CARACTERISTICA} permite una flexibilidad clave-valor (Entity-Attribute-Value), pero esto conlleva el riesgo de datos sucios (ej. "Rojo", "rojo ", "Color Rojo"). Este disparador normaliza los datos y restringe los valores permitidos según el tipo de característica.

\textbf{Descripción:} Realiza dos funciones:
\begin{enumerate}
    \item \textbf{Normalización:} Convierte claves y valores a mayúsculas y elimina espacios en blanco.
    \item \textbf{Validación Lógica:} Utiliza una estructura \texttt{CASE} para aplicar listas blancas de valores permitidos según la etiqueta (Colores válidos, Tamaños válidos, Prioridades).
\end{enumerate}

\begin{lstlisting}[language=SQL, caption={Trigger TRG\_VALIDAR\_COHERENCIA}]
CREATE OR REPLACE TRIGGER TRG_VALIDAR_COHERENCIA
BEFORE INSERT OR UPDATE ON CARACTERISTICA
FOR EACH ROW
DECLARE
    v_nombre_norm VARCHAR2(100);
    v_valor_norm VARCHAR2(100);
BEGIN
    -- 1. Normalizamos los datos (quitamos espacios y pasamos a mayúsculas)
    v_nombre_norm := UPPER(TRIM(:NEW.NOMBRE));
    v_valor_norm := UPPER(TRIM(:NEW.VALOR));

    -- 2. Lógica de control (Switch case)
    CASE v_nombre_norm
        -- CASO 1: Validación de Colores
        WHEN 'COLOR' THEN
            IF v_valor_norm NOT IN ('ROJO', 'AZUL', 'VERDE', 'AMARILLO', 'NEGRO', 'BLANCO', 'GRIS') THEN
                RAISE_APPLICATION_ERROR(-20005,
                'Valor incoherente: "' || :NEW.VALOR || '" no es un color válido.');
            END IF;

        -- CASO 2: Validación de Tamaño
        WHEN 'TAMANO' THEN
            IF v_valor_norm NOT IN ('PEQUEÑO', 'MEDIANO', 'GRANDE', 'BANNER', 'FULLSCREEN') THEN
                RAISE_APPLICATION_ERROR(-20005,
                'Valor incoherente: "' || :NEW.VALOR || '" no es un tamaño válido.');
            END IF;

        -- CASO 3: Validación de Prioridad
        WHEN 'PRIORIDAD' THEN
            IF v_valor_norm NOT IN ('ALTA', 'MEDIA', 'BAJA') THEN
                RAISE_APPLICATION_ERROR(-20005, 
                'La prioridad solo puede ser Alta, Media o Baja.');
            END IF;

        ELSE
            NULL; -- Se permiten otras características sin restricción
    END CASE;
END;
/
\end{lstlisting}